\subsubsection{Evaluation on the Fock-module ends}
\label{mbh:sec}

The derived tensor product in Theorem~\ref{mb:thm:mode-bar} can be
evaluated without changing either module end.  Its differential reduces
to the right action of the annihilation operators on the restricted
dual.  An explicit comparison works over an arbitrary commutative base;
the homology of the resulting complex then depends on the level and
the characteristic.

For the comparison, let \(S\) be a commutative ring and fix
\(c_m\in S\), \(m\ge1\).  Let \(W\) be the unital \(S\)-algebra
with relations
\[
 [p_m,q_n]=c_m\delta_{mn},\qquad
 [p_m,p_n]=[q_m,q_n]=0,
\]
and put \(V_S=S[q_1,q_2,\ldots]\), with
\(q_m=q_m\cdot\) and \(p_m=c_m\partial_{q_m}\).
The preceding mode algebra is obtained by taking \(S=k\) and
\(c_m=\kappa m\), with \(q_m=h_{m-1}\).
All words, polynomials, exterior tensors and sums here have finite
algebraic support.

The ordered-monomial argument above is valid over \(S\).
Indeed, \(q_m=Q_m\cdot\) and
\(p_m=c_m\partial_{Q_m}+P_m\cdot\) act on
\(S[Q_m,P_m:m\ge1]\).  The represented ordered monomial
\(q^\alpha p^\beta\) sends \(1\) to \(Q^\alpha P^\beta\).
Together with reordering, this proves that
the \(q^\alpha p^\beta\) form an \(S\)-basis of \(W\), including
when some \(c_m\) vanish.

Let \(E_S=\bigoplus_{m\ge1}S\epsilon_m\), with each
\(\epsilon_m\) in homological degree one.  For a finite ordered set
\(I=\{i_1<\cdots<i_n\}\), write
\(\epsilon_I=\epsilon_{i_1}\wedge\cdots\wedge\epsilon_{i_n}\).
Define a complex of free left \(W\)-modules by
\begin{equation}\label{mbh:eq:free-koszul}
 \begin{aligned}
 K^W_n&=W\otimes_S\Lambda^n E_S,\\
 d_K(w\epsilon_I)
   &=\sum_{a=1}^n(-1)^{a-1}
       wp_{i_a}\epsilon_{I\setminus\{i_a\}},&
 \pi_K(w)&=w1\in V_S.
 \end{aligned}
\end{equation}
Commutation of the \(p_m\) gives \(d_K^2=0\).
For comparison, write
\[
 P^W_n=W\otimes_S W^{\otimes_S n}\otimes_S V_S
\]
for the unnormalized left bar resolution.  Its differential \(d_P\)
has the faces of Theorem~\ref{mb:thm:mode-bar}, with the first
coefficient in the regular left module \(W\), and its augmentation is
\(\pi_P(a[]v)=av\).  It is free as a left \(W\)-module on the basis
\begin{equation}\label{mbh:eq:bar-basis}
 e=1[a_1|\cdots|a_n]q^\alpha,
\end{equation}
where every \(a_j\) is an ordered monomial.

\begin{lemma}[A contraction over the coefficient ring]
\label{mbh:lem:integral-contraction}
For an ordered basis vector \(q^\alpha p^\beta\epsilon_I\),
put \(J=\operatorname{supp}\beta\cup I\).
If \(J=\varnothing\), set \(h_K(q^\alpha)=0\).
Otherwise let \(m=\min J\) and set
\begin{equation}\label{mbh:eq:least-index}
 h_K(q^\alpha p^\beta\epsilon_I)=
 \begin{cases}
 0,&m\in I,\\
 q^\alpha p^{\beta-e_m}\epsilon_m\wedge\epsilon_I,&m\notin I.
 \end{cases}
\end{equation}
Here \(e_m\) is the multi-index with \(1\) in position \(m\)
and \(0\) elsewhere.  Extend this map by \(S\)-linearity, and let
\(i_K:V_S\to K^W_0\) be the polynomial inclusion.  Then
\begin{equation}\label{mbh:eq:koszul-contraction}
 d_Kh_K+h_Kd_K=1-i_K\pi_K,\qquad h_K^2=0,
\end{equation}
where \(i_K\pi_K\) acts only in degree zero.  In particular,
\(K^W_\bullet\to V_S\) is a free resolution.
\end{lemma}

\begin{proof}
If \(J=\varnothing\), the input is a polynomial in the \(q_m\)
in degree zero, and both sides of the first identity vanish.
If \(m\in I\), the homotopy vanishes on the input.  In its
differential, only the term removing \(\epsilon_m\) has nonzero
image under \(h_K\).  Since \(m\) is the first exterior index,
this term has positive sign and recovers the input.

If \(m\notin I\), then \(\beta_m>0\) and \(m<i\) for every
\(i\in I\).  Thus the exterior word in
\eqref{mbh:eq:least-index} is already ordered.  The differential
term removing its initial \(\epsilon_m\) recovers the input.
Every other term cancels the term obtained by first removing the
same exterior generator and then applying \(h_K\): the insertion
of the initial generator changes its sign once.  This proves the
first identity.  After a nonzero application of \(h_K\), the least
occupied index belongs to the exterior set, so a second application
is zero.  Taking \(h_{K,-1}=i_K\) contracts the augmented complex.
\end{proof}

The bar contraction from the proof of
Theorem~\ref{mb:thm:mode-bar} also works over \(S\):
\begin{equation}\label{mbh:eq:bar-contraction}
 h_{P,-1}(v)=1[]v,\qquad
 h_{P,n}\bigl(a_0[a_1|\cdots|a_n]v\bigr)
       =1[a_0|a_1|\cdots|a_n]v.
\end{equation}
Its first face is the input and the other faces cancel those in
\(h_Pd_P\), so \(d_Ph_P+h_Pd_P=1\) on the augmented complex.
Both \(h_K\) and \(h_P\) are only \(S\)-linear.
For example, \(h_K(p_m)=\epsilon_m\), whereas
\(p_mh_K(1)=0\).  The module-linear comparison must therefore be
defined on free bases and then extended by \(W\)-linearity.

Define \(F:K^W\to P^W\) by antisymmetrization:
\begin{equation}\label{mbh:eq:antisymmetrization}
 F_n(w\epsilon_I)=
 \sum_{\sigma\in\mathfrak S_n}\operatorname{sgn}(\sigma)\,
 w[p_{i_{\sigma(1)}}|\cdots|p_{i_{\sigma(n)}}]1,
 \qquad F_0(w)=w[]1.
\end{equation}
Here \(\mathfrak S_n\) is the permutation group of \(\{1,\ldots,n\}\).
For the reverse map, prescribe
\(G_0(1[]q^\alpha)=q^\alpha\), and, for \(n>0\), prescribe
on each basis element \(e\) of \eqref{mbh:eq:bar-basis}
\begin{equation}\label{mbh:eq:reverse-recursion}
 G_n(e)=h_{K,n-1}\bigl(G_{n-1}(d_Pe)\bigr).
\end{equation}
Expand the argument in ordered monomials before applying \(h_K\),
and extend each \(G_n\) by left \(W\)-linearity.
Finally put \(D=1_{P^W}-FG\), set \(H_{-1}=0\), and define
on the same free basis
\begin{equation}\label{mbh:eq:bar-homotopy-recursion}
 H_n(e)=h_{P,n}\bigl(D_n(e)-H_{n-1}d_Pe\bigr),
 \qquad n\ge0,
\end{equation}
with the second term omitted for \(n=0\).
Again extend by left \(W\)-linearity.  These prescriptions apply
to the displayed free basis; they do not apply the \(S\)-linear
contractions to an arbitrary left multiple of a basis element.

\begin{theorem}[An integral retraction of the oscillator bar]
\label{mbh:thm:retraction}
The maps \(F\) and \(G\) are left \(W\)-linear chain maps lifting
the identity of \(V_S\), and
\begin{equation}\label{mbh:eq:retraction-identities}
 GF=1_{K^W},\qquad d_PH+Hd_P=1_{P^W}-FG.
\end{equation}
All formulas terminate on each element, introduce no new oscillator
index, and commute with every base change \(S\to S'\), with each
\(c_m\) replaced by its image in \(S'\).  They use no division.
\end{theorem}

\begin{proof}
For \(F\), the first bar faces give \(Fd_K\).
The last faces vanish since \(p_m1=0\), and the interior faces
cancel in pairs by interchanging two adjacent commuting \(p_m\).
The signs in each pair are opposite, also in characteristic two.
Thus \(d_PF=Fd_K\), and the augmentations agree.

The identity \(\pi_KG_0=\pi_P\) follows on the free basis.
Assume the chain identity for \(G\) in lower degrees, and let
\(e\in P^W_n\) be a free basis element.  Put
\(z=G_{n-1}d_Pe\).  If \(n=1\), then
\(\pi_Kz=\pi_Pd_Pe=0\).  If \(n>1\), then
\(d_Kz=G_{n-2}d_P^2e=0\).  In either case,
\eqref{mbh:eq:koszul-contraction} gives \(d_Kh_Kz=z\).
This proves \(d_KG_n(e)=G_{n-1}d_Pe\).
Both sides are \(W\)-linear, so the identity holds everywhere.

In degree zero, \(G_0F_0=1\).  Suppose \(GF=1\) in degree
\(n-1\), and take \(n>0\).
The vector \(F_n(\epsilon_I)\) is an \(S\)-linear combination
of free basis elements of \(P^W_n\).  Applying
\eqref{mbh:eq:reverse-recursion} term by term therefore gives
\[
 G_nF_n(\epsilon_I)
   =h_KG_{n-1}d_PF_n(\epsilon_I)
   =h_Kd_K(\epsilon_I)=\epsilon_I.
\]
For the last equality, \(h_K(\epsilon_I)=0\) and
\eqref{mbh:eq:koszul-contraction} applies in positive degree.
Left \(W\)-linearity proves the first identity in
\eqref{mbh:eq:retraction-identities}.

For the second identity, \(\pi_PD_0=0\), since \(F\) and \(G\)
lift the identity of \(V_S\).  Thus
\(d_PH_0(e)=D_0(e)\) by \eqref{mbh:eq:bar-contraction}.
Assume the homotopy identity below degree \(n\), and put
\(z=D_n(e)-H_{n-1}d_Pe\).
Since \(D\) is a chain map, the preceding degree gives
\[
 d_Pz
  =D_{n-1}d_Pe
    -(D_{n-1}-H_{n-2}d_P)d_Pe=0.
\]
The bar contraction yields \(d_Ph_Pz=z\), proving the required
identity on the free basis and hence on all left multiples.

Each recursion lowers homological degree.  Reordering a finite
product gives finitely many ordered monomials, and none of the
contractions introduces an index absent from the input.  Thus the
maps are compatible with inclusions of finite sets of modes, with
their order inherited from the positive integers.
Their coefficients use only addition, multiplication and signs in
the \(c_m\).  The same free-basis formulas after any scalar change
satisfy the same identities; no flatness assumption on that change
is required.
\end{proof}

In particular, antisymmetrization does not require division by a
factorial.  For \(i<j\), the recursion gives
\[
 \begin{aligned}
 G_1(1[p_i]1)&=\epsilon_i,&
 G_1(1[p_ip_j]1)&=p_j\epsilon_i,\\
 G_2(1[p_i|p_j]1)&=\epsilon_i\wedge\epsilon_j,&
 G_2(1[p_j|p_i]1)&=0.
 \end{aligned}
\]
Their antisymmetric difference maps to
\(\epsilon_i\wedge\epsilon_j\), including in characteristic two.

For any right \(W\)-module \(L\), tensoring the maps and homotopy
over \(W\) gives a chain homotopy equivalence
\begin{equation}\label{mbh:eq:module-comparison}
 B_S(L,W,V_S)
 \ \underset{F_L}{\stackrel{G_L}{\rightleftarrows}}\
 K(L),\qquad
 K(L)_n=L\otimes_S\Lambda^n E_S,
\end{equation}
where \(F_L=1_L\otimes_WF\) goes from \(K(L)\) to the bar,
\(G_L=1_L\otimes_WG\) goes from the bar to \(K(L)\), and
\begin{equation}\label{mbh:eq:annihilator-differential}
 d(\lambda\epsilon_I)=
 \sum_{a=1}^n(-1)^{a-1}
       (\lambda p_{i_a})\epsilon_{I\setminus\{i_a\}}.
\end{equation}
The identities \(G_LF_L=1\) and
\(dH_L+H_Ld=1-F_LG_L\), with \(H_L=1_L\otimes_WH\),
follow by tensoring \eqref{mbh:eq:retraction-identities}.
This requires no flatness of \(L\) over \(W\).

We now use the actual restricted-dual end.  Give \(q_m\) weight
\(m\), write \((V_S)_a\) for the weight-\(a\) component, and put
\[
 M_S=\bigoplus_{a\ge0}\operatorname{Hom}_S((V_S)_a,S).
\]
Each weight component is finite free.  Its monomial duals form a
basis \(\lambda_\alpha\), characterized by
\(\lambda_\alpha(q^\beta)=\delta_{\alpha,\beta}\).
The right action \((\lambda a)(v)=\lambda(av)\) gives
\begin{equation}\label{mbh:eq:dual-actions}
 \lambda_\alpha p_m
     =c_m(\alpha_m+1)\lambda_{\alpha+e_m},
 \qquad
 \lambda_\alpha q_m=
 \begin{cases}
 \lambda_{\alpha-e_m},&\alpha_m>0,\\
 0,&\alpha_m=0.
 \end{cases}
\end{equation}
Thus \eqref{mbh:eq:annihilator-differential} computes the
derived tensor product on the stated Fock-module ends over \(S\).
The evaluation of the bar, \(\lambda[]v\mapsto\lambda(v)\)
in degree zero, becomes
\begin{equation}\label{mbh:eq:evaluation}
 \pi(\lambda_\alpha)=\delta_{\alpha,0}
\end{equation}
on \(K(M_S)_0\), with zero value in positive degrees.
For \(F_0\) this follows from \(F_0(w)=w[]1\).
For \(G_0\) it follows from
\((\lambda aq^\alpha)(1)=\lambda(aq^\alpha1)\).
The map is a chain map because \((\lambda p_m)(1)=0\).

\begin{theorem}[Evaluation at nonzero level]
\label{mbh:thm:nonzero}
Let \(k\) be a field of characteristic zero and \(c\in k^\times\).
For \(\mathscr W_c\), \(V=k[q_1,q_2,\ldots]\) and its
restricted dual \(M\), the evaluation
\[
 B_k(M,\mathscr W_c,V)\longrightarrow k
\]
is a quasi-isomorphism.  Equivalently,
\[
 \operatorname{Tor}^{\mathscr W_c}_j(M,V)=
 \begin{cases}k,&j=0,\\0,&j>0.\end{cases}
\]
\end{theorem}

\begin{proof}
All factors in
\begin{equation}\label{mbh:eq:fixed-normalization}
 k[x_1,x_2,\ldots]\longrightarrow M,\qquad
 x^\alpha\longmapsto
 \left(\prod_m(cm)^{\alpha_m}\alpha_m!\right)\lambda_\alpha
\end{equation}
are invertible.  This is a vector-space isomorphism, and
\eqref{mbh:eq:dual-actions} identifies right multiplication by
\(p_m\) with multiplication by \(x_m\).
The complex \(K(M)\) becomes
\[
 K=k[x_1,x_2,\ldots]\otimes_k\Lambda(\epsilon_m:m\ge1),
 \qquad
 \delta=\sum_m x_m\iota_m,
\]
where \(\iota_m\) is exterior contraction:
\(\iota_m(\epsilon_I)=(-1)^{a-1}
\epsilon_{I\setminus\{m\}}\) when \(m=i_a\), and zero
when \(m\notin I\).

Give both \(x_m\) and \(\epsilon_m\) auxiliary weight one.
This is separate from the state weight and homological degree.
On auxiliary weight \(w>0\), define
\begin{equation}\label{mbh:eq:euler-homotopy}
 h=\frac1w\sum_m\epsilon_m\wedge\partial_{x_m},
 \qquad h(1)=0.
\end{equation}
The relations
\(\iota_i(\epsilon_j\wedge-)+\epsilon_j\wedge\iota_i
=\delta_{ij}\) and the polynomial product rule give
\[
 \delta\Bigl(\sum_m\epsilon_m\wedge\partial_{x_m}\Bigr)
 +\Bigl(\sum_m\epsilon_m\wedge\partial_{x_m}\Bigr)\delta
 =\sum_m x_m\partial_{x_m}+\sum_m\epsilon_m\wedge\iota_m.
\]
The last operator is multiplication by \(w\).
Thus \(\delta h+h\delta=1-i\pi\), where \(i(1)=1\) and
\(\pi\) is the constant coefficient in degree zero.
Every sum is finite on its input.  This contracts \(K\) onto
\(k\), with exactly the evaluation \eqref{mbh:eq:evaluation}.
The equivalence \eqref{mbh:eq:module-comparison} proves the claim.
\end{proof}

This scalar homology uses the two module actions.
It supplies no character of \(\mathscr W_c\), which remains
excluded by \([p_1,q_1]=c1\).
Over a parameter ring, division by the level would lose homology.

\begin{theorem}[The level family and derived specialization]
\label{mbh:thm:family}
Let \(k\) be a field of characteristic zero, put \(S=k[\kappa]\),
and take \(c_m=\kappa m\).  The bar on \(V_S\) and \(M_S\)
is chain homotopy equivalent to
\begin{equation}\label{mbh:eq:family-complex}
 K_S=S[x_1,x_2,\ldots]\otimes_S\Lambda(\epsilon_m:m\ge1),
 \qquad d_\kappa=\kappa\delta,\qquad
 \delta=\sum_m x_m\iota_m.
\end{equation}
Set \(B_j=\operatorname{im}(\delta:K_{S,j+1}\to K_{S,j})\)
and \(B_{-1}=0\).  Then
\begin{equation}\label{mbh:eq:family-homology}
 H_j(K_S)=
 \begin{cases}
 S\oplus B_0/\kappa B_0,&j=0,\\
 B_j/\kappa B_j,&j>0,
 \end{cases}
 \qquad
 B_0=(x_1,x_2,\ldots)S[x_1,x_2,\ldots].
\end{equation}
The kernel of evaluation on homology,
\(H_\bullet(K_S)\to S\) with \(S\) in degree zero, is
annihilated by \(\kappa\).  It is nonzero in every positive
homological degree.
The derived zero fibre is
\begin{equation}\label{mbh:eq:derived-zero}
 H_j(K_S\otimes_S^{\mathbb L}k)
   =k[x_1,x_2,\ldots]\otimes_k\Lambda^j(\epsilon_m:m\ge1),
 \qquad k=S/(\kappa).
\end{equation}
\end{theorem}

\begin{proof}
Normalize the monomial duals by
\[
 x^\alpha\longmapsto
 \left(\prod_m m^{\alpha_m}\alpha_m!\right)\lambda_\alpha.
\]
The factors are units in \(k\), so this is an invertible
\(S\)-linear change of basis.  It introduces no inverse of
\(\kappa\).  Formula~\eqref{mbh:eq:dual-actions} then gives
\eqref{mbh:eq:family-complex}.
The homotopy \eqref{mbh:eq:euler-homotopy} is \(S\)-linear and
satisfies
\begin{equation}\label{mbh:eq:family-homotopy}
 d_\kappa h+hd_\kappa=\kappa(1-i\pi).
\end{equation}
Here \(\pi\) takes the constant coefficient in the variables
\(x_m\), with values in \(S\), and \(i\) is the scalar inclusion.

Let \(K_{S,j}^+\) denote the part of positive auxiliary weight;
only the scalar summand \(S\cdot1\) in degree zero is removed.
The identity \(\delta h+h\delta=1\) on this part gives
\begin{equation}\label{mbh:eq:family-splitting}
 K_{S,j}^+=B_j\oplus hB_{j-1},\qquad
 d_\kappa(hb)=\kappa b,\qquad d_\kappa(B_j)=0.
\end{equation}
Indeed, every \(z\) is
\(\delta hz+h\delta z\), which lies in the stated sum.
If \(hb\in B_j\) with \(b\in B_{j-1}\), applying
\(\delta\) gives \(b=0\), proving directness.
Since \(\delta\) and \(h\) are defined over \(k\), each \(B_j\)
is \(S\) tensored with a \(k\)-vector space and hence is
\(S\)-free.  Multiplication by \(\kappa\) is injective on it.
The kernels and images in \eqref{mbh:eq:family-splitting}
now give \eqref{mbh:eq:family-homology}.

For \(j>0\), the cycle
\begin{equation}\label{mbh:eq:torsion-cycle}
 z_j=\delta(\epsilon_1\wedge\cdots\wedge\epsilon_{j+1})
   =\sum_{a=1}^{j+1}(-1)^{a-1}x_a\,
      \epsilon_1\wedge\cdots\wedge
      \widehat{\epsilon_a}\wedge\cdots\wedge\epsilon_{j+1}
\end{equation}
lies in \(B_j\) and has nonzero coefficients independent of
\(\kappa\).  It cannot lie in \(\kappa B_j\), so its class is
nonzero.  The equality
\(d_\kappa(\epsilon_1\wedge\cdots\wedge\epsilon_{j+1})
=\kappa z_j\) also displays its annihilation.

The splitting \eqref{mbh:eq:family-splitting} expresses \(K_S\)
as the scalar \(S\) and a direct sum of two-term free complexes
\[
 hB_j\xrightarrow{\ \kappa\ }B_j,
\]
in degrees \(j+1\) and \(j\).  A two-term complex of free
modules is \(K\)-flat: tensoring it with an acyclic complex is
the cone of a map between acyclic direct sums of that complex.
Direct sums preserve exactness.  The same is therefore true of
\(K_S\), so its ordinary tensor product with \(k=S/(\kappa)\)
computes the derived tensor product.  Its differential becomes zero,
proving \eqref{mbh:eq:derived-zero}.
\end{proof}

The extra specialization classes can be identified without a
convergence argument.  For \(j\ge1\), the preceding splitting gives
the exact sequence
\begin{equation}\label{mbh:eq:base-change}
 \begin{aligned}
 0\longrightarrow H_j(K_S)\otimes_S k
 &\longrightarrow H_j(K_S\otimes_S^{\mathbb L}k)\\
 &\longrightarrow
 \operatorname{Tor}_1^S(H_{j-1}(K_S),k)
 \longrightarrow0.
 \end{aligned}
\end{equation}
Its first term is represented by \(B_j\otimes_Sk\), and its
last term by \(hB_{j-1}\otimes_Sk\), using
\(h:B_{j-1}\to hB_{j-1}\).
Indeed, \(0\to B_{j-1}\xrightarrow{\kappa}B_{j-1}
\to B_{j-1}/\kappa B_{j-1}\to0\) is a free resolution,
and the scalar summand of \(H_0\) has zero \(\operatorname{Tor}_1\).
For a single oscillator, \(B_1=0\): the family has \(H_1=0\),
whereas its zero fibre has
\(H_1=k[x_1]\epsilon_1\).  These classes come entirely from
the last term of \eqref{mbh:eq:base-change}.
With infinitely many oscillators,
\eqref{mbh:eq:torsion-cycle} gives family torsion in every
positive degree.  Inverting \(\kappa\) instead gives the
contraction \(\kappa^{-1}h\) and recovers
Theorem~\ref{mbh:thm:nonzero}.

At zero level over any field, every right \(p_m\)-action on
the restricted dual vanishes.  Thus
\eqref{mbh:eq:annihilator-differential} gives
\begin{equation}\label{mbh:eq:zero-level}
 \operatorname{Tor}^{\mathscr W_0}_j(M,V)
   =M\otimes_k\Lambda^j(\epsilon_m:m\ge1).
\end{equation}
The algebra \(\mathscr W_0\) still contains every \(p_m\);
its zero action on these module ends does not remove the
corresponding exterior generator.

\begin{theorem}[The positive-characteristic basis]
\label{mbh:thm:positive-characteristic}
Let \(k\) be a field of characteristic \(\ell>0\), let \(c\in k^\times\),
and use the polynomial Fock module and its restricted dual.
Put
\[
 G=\{m\ge1:\ell\nmid m\},\qquad
 Z=\{m\ge1:\ell\mid m\}.
\]
For every \(j\ge0\), a basis of
\(\operatorname{Tor}^{\mathscr W_c}_j(M,V)\) is represented
in \(K(M)\) by
\begin{equation}\label{mbh:eq:characteristic-basis}
 \lambda_{\alpha(b,I)}\epsilon_I,\qquad
 |I|=j,\qquad
 \alpha(b,I)_m=
 \begin{cases}
 \ell b_m+(\ell-1)\mathbf1_{m\in I},&m\in G,\\
 b_m,&m\in Z,
 \end{cases}
\end{equation}
where \(I\) is a finite set and \(b=(b_m)\) is a finite-support
multi-index of nonnegative integers.
This is a basis of graded vector spaces; no multiplication on
these Tor groups is asserted.
\end{theorem}

\begin{proof}
For a single mode \(m\), let
\(M_m=\bigoplus_{n\ge0}k\lambda_n\).
Its factor in the annihilator complex is
\begin{equation}\label{mbh:eq:one-mode}
 M_m\epsilon_m\longrightarrow M_m,\qquad
 \lambda_n\epsilon_m\longmapsto cm(n+1)\lambda_{n+1}.
\end{equation}
If \(m\in G\), its degree-zero cokernel has basis
\(\lambda_{\ell b}\), \(b\ge0\), and its degree-one kernel
has basis \(\lambda_{\ell b+\ell-1}\epsilon_m\).
Every other degree-zero basis vector is the image of exactly one
basis vector with nonzero coefficient, and every remaining
degree-one vector belongs to such a pair.
If \(m\in Z\), the differential is zero and both copies of
\(M_m\) survive.

The full complex is the algebraic restricted tensor product of
\eqref{mbh:eq:one-mode} over all \(m\), using \(\lambda_0\)
in degree zero at each new index.  For a finite set of indices,
split each factor into the displayed kernel and cokernel
representatives and its two-term pairs with nonzero differential.
Those pairs are contractible over the field \(k\).  Tensoring their
contractions with the other factors proves that the homology of the
finite tensor product is the tensor product of the displayed
homologies.  This splitting preserves the distinguished
degree-zero vector \(\lambda_0\), so it is compatible with
inclusions of finite index sets.  Their union is the full
algebraic complex, and filtered unions of vector spaces preserve
exactness.  Selecting the degree-one factor exactly at the indices
of \(I\) gives precisely \eqref{mbh:eq:characteristic-basis}.
\end{proof}

For example,
\[
 [\,\lambda_{\ell-1}[p_1]1\,]
   \in H_1\bigl(B_k(M,\mathscr W_c,V)\bigr)
\]
is nonzero.  Both bar faces vanish: the first has coefficient
\(c\ell=0\), and the last contains \(p_1 1=0\).
Under \(F_M\), it is the image of
\(\lambda_{\ell-1}\epsilon_1\), whose class is nonzero by
Theorem~\ref{mbh:thm:positive-characteristic};
\(G_MF_M=1\) proves that it is not a boundary.
Thus the evaluation to \(k\) fails to be a quasi-isomorphism
in every positive characteristic, including at nonzero level.
Modes in \(Z\) supply further independent classes.
Changing the polynomial Fock module to a divided-power module
would change the input action in \eqref{mbh:eq:dual-actions}.

These evaluations retain the two module ends and concern ordinary
algebraic complexes.  The balanced cuts of
\eqref{mb:eq:module-cuts} retain their stated target; the
retraction \eqref{mbh:eq:module-comparison} makes no assertion
of a coproduct equivalence.  A geometric collision comparison
or an extension to completed states requires additional maps.
